% Section 1: Introduction
% NOTE: \section{Introduction} and \label{sec:intro} are in main.tex

A physician documents ``patient denies chest pain, sister had MI at 45, will consider statin if lipids remain elevated''---a single sentence encoding four epistemic states: a negated symptom, a family-history event, a hypothetical action, and an implicit present condition.
Clinical NLP can detect these assertions accurately~\cite{uzuner2011i2b2assertion, gul2025beyondnegation}, yet when a RAG system ingests the same note, assertion context is flattened: negation is lost, family attribution collapses onto the patient, and downstream reasoning may conclude the patient \emph{has} chest pain and \emph{had} an MI.
We call this the \emph{epistemic propagation gap}---the systematic loss of assertion status as clinical information moves from extraction through storage to retrieval.

The gap persists because, to our knowledge, prior KG-RAG systems do not propagate assertion labels beyond the extraction stage into graph-augmented retrieval.
Clinical NLP pipelines (cTAKES, John Snow Labs) detect assertions, but do not carry them as edge attributes into knowledge graph construction or retrieval.
OMOP excludes negated conditions from \texttt{CONDITION\_OCCURRENCE}~\cite{omop2024}; FHIR provides \texttt{verificationStatus} for \texttt{Condition} resources only, covering a fraction of the clinical data model.
Graph-augmented RAG systems---GraphRAG~\cite{edge2024graphrag}, GFM-RAG~\cite{luo2025gfmrag}, KARE~\cite{jiang2025kare}, Medical-Graph-RAG~\cite{wu2025medicalgraphrag}---build KGs that discard the metadata distinguishing ``patient has diabetes'' from ``rule out diabetes.''
Similarly, recent clinical QA systems---MedPaLM~2~\cite{singhal2023medpalm2}, Med-Gemini~\cite{saab2024medgemini}, AMIE~\cite{tu2024amie}---excel at medical knowledge recall but are not designed for longitudinal assertion reasoning over real patient records, which is the gap ClinicalBench targets.
A parallel \emph{temporal integration gap} exists: clinical events span three time dimensions (valid time, transaction time, NLP-asserted temporality), yet existing systems model at most two~\cite{huang2024medtkg,ramirez2025graphiti}.

We present \epikg, a clinical KG-RAG system that preserves epistemic status end-to-end, and ClinicalBench, a 400-question benchmark that exposes the propagation gap.
Our contributions are primarily \emph{empirical findings and a design principle}, not system components:

\begin{enumerate}[leftmargin=*, topsep=4pt, itemsep=2pt]
  \item \textbf{The epistemic propagation gap.} We identify and formalize the systematic loss of assertion metadata across clinical NLP pipelines, derive an information-theoretic loss bound (Section~\ref{sec:system:formal}, Appendix~\ref{app:formal}), and implement end-to-end preservation as a testable invariant.

  \item \textbf{Assertions alone do not help; routing is the active ingredient.} Adding assertion metadata to graph edges without question-type-specific retrieval does not improve accuracy ($-3.2\pp$, C3$\to$C4; CI crossing zero, $p = 0.37$). The entire C3$\to$C4g gain ($+19.0\pp$) is attributable to intent-aware routing ($+22.2\pp$, C4$\to$C4g; $p < 10^{-10}$). Metadata enrichment without retrieval alignment is insufficient.

  \item \textbf{Knowledge-compensator pattern.} Across four models from distinct training lineages, models with weaker baselines gain more from KG-RAG, and all converge to 57.5--69.0\% under C4g despite C1 baselines spanning 20.5\% to 37.0\%. This convergence pattern---observed descriptively across four data points---is consistent with the epistemic KG providing a knowledge floor that partially compensates for model-specific weaknesses.

  \item \textbf{ClinicalBench.} A 400-question benchmark across 9 assertion-sensitive categories (negation, uncertainty, conditional, family history, sequence, current state, duration, historical, change) that exposes category$\times$condition interactions invisible to aggregate metrics. Released with Croissant metadata on HuggingFace.

  \item \textbf{A falsifiable design principle.} Preserve epistemic metadata \emph{and} route retrieval by question intent, then evaluate by category$\times$condition interaction rather than aggregate score. This is not a system architecture but a falsifiable constraint: any pipeline violating it should underperform on assertion-sensitive longitudinal QA.
\end{enumerate}

Our central research question is: \emph{under what conditions does structured epistemic context help, hurt, or break even relative to unstructured retrieval and no retrieval?}
The answer is interactional rather than uniform: gains depend on question type, routing policy, and model.
Intent-aware KG-RAG improves all nine ClinicalBench categories, with largest effects on cross-admission longitudinal tasks; assertions without routing do not improve temporal categories; and brute-force long context does not recover these failures despite a large token window.
